% method.tex -- source: research/paper2/method/GATE-METHOD-SPEC-v2-2026-10-28.md (v1 where v2 says unchanged).
% The only empirical numbers in this file are the M1 detection values quoted in Section "gate statistic" (run P2-GATE-M1, K2 wording).
% The only guarantee stated is the one the SAP (section 6) allows.
% Notation: symbols follow the SAP and exp_p2/gate/theory.py. Four letters with a second meaning were disambiguated in this revision
% (groups Gamma_k, sample count n_draw, empirical calibration risk F_cal, SNR flag probabilities phi_0, phi_1); crosswalk in Appendix A.
% Propositions: Prop. 1 (calibration) is conformal risk control applied, not a contribution; the marginalisation property is
% the standard denoiser property, not a contribution (must-not-claim items 5 and 6).
% Prop. 2 (correlated-token gating) is theory (D0078, D0079); it is never used to explain an E-pred result; Props. 1 and 3 of the theory addendum are not used here.

This section defines the notation, the scores, the calibrated gate, the evidence statistic and seven algorithms.
Each algorithm calls the earlier ones, and Algorithm~\ref{alg:gate} composes all of them.
The method has three steps.
Each input token is \emph{scored} by its masked log-likelihood inside its own modality with the other modality withheld.
Tokens whose score is improbable for clean inputs, at a threshold set on clean data, are \emph{gated}.
The decision then \emph{fuses} evidence from the tokens that remain (GATE-M, the marginalise arm), or from tokens re-infilled by
the model (GATE-I, the impute arm; the SAP calls it I-GATE).
Cost is counted in network passes, listed per quantity in \ref{app:diag}.
One design rule organises the three steps: evidence that is tested must not be used to gate (Section~\ref{ssec:separation}).
The denoiser is used through one conditional only, the probability of a token given a visible subset, called with no time
input; Section~\ref{ssec:data} and \ref{app:timefree} state what is checked about that property for the released
checkpoint.

\subsection{Setting and notation}\label{ssec:notation}
A unified masked discrete diffusion model with absorbing-state masking is given a sequence of caption tokens and image tokens.
The network $p_\theta$ maps a partly masked sequence to a categorical distribution at every masked position.
A pair is $x=(t,v)$ with caption tokens $t=(t_1,\dots,t_n)$ and image tokens $v=(v_1,\dots,v_m)$.
The modalities are $a,b\in\{T,I\}$ with $a\neq b$.
Table~\ref{tab:notation} lists the symbols.
All log-likelihoods are in nats, computed as an \fp{} log-softmax of the logits restricted to the vocabulary of the scored
modality.

\begin{table}[t]
\centering
\caption{Notation. Symbols follow the preregistration documents and the analysis code. A letter with a second use appears in a
separate row. Symbols used only in the protocol are defined there.}
\label{tab:notation}
\footnotesize
\begin{tabular}{p{0.27\linewidth}p{0.63\linewidth}}
\toprule
Symbol & Meaning \\
\midrule
$x=(t,v)$ & one pair of caption tokens $t$ ($n$ tokens) and image tokens $v$ ($m=1\,024$ tokens) \\
$\tilde x=(\tilde t,\tilde v)$, $v^\circ$ & the observed, possibly corrupted pair; the tokens of the clean image \\
$Z_j=\ind[\tilde v_j\neq v^\circ_j]$ & corrupted-token label (used only for evaluation) \\
$\mask$, $\varnothing^a$ & the mask token ($\mask$ means ``unknown''); all tokens of modality $a$ masked \\
$a,b\in\{T,I\}$ & scored modality and the other modality (text, image) \\
$p_\theta(x_j=\cdot\mid\cdot)$ & categorical of the network at masked position $j$ \\
$\Pi_K=\{\Gamma_1,\dots,\Gamma_K\}$ & random partition of the positions of modality $a$ into $K$ complementary groups; $k(j)$ is the group of $j$ \\
$A_j,B_j,C_j,D_j$ & four context levels of Eq.~\eqref{eq:levels}; $D_j[y]$ is the vector over codes $y$ and $D_j=D_j[x_j]$ \\
$q_j$, $u_j$ & hypothesis-free intra-modal evidence $B_j-D_j$ and suspicion $u_j=-q_j$ \\
$\alpha$, $\hat\tau$, $\widehat F_{\rm cal}$ & target false-alarm level, calibrated threshold, empirical calibration risk (Eq.~\ref{eq:crc}) \\
$G$, $R$ & gated (flagged) set and reliable set, the complement of $G$ \\
$c$, $c^+$, $c^-$ & a candidate caption, the true caption and a wrong caption \\
$\Lambda^c_j[y]$ & log-probability of image code $y$ at position $j$ with the image fully masked and caption $c$ visible \\
$\Delta_j$, $\bar\Delta_R$, $Y_R$ & per-token paired evidence, its mean over $R$, and the decision $\ind[\bar\Delta_R>0]$ (Eq.~\ref{eq:delta}) \\
$d_R$, $z_R$ & size-normalised evidence ratio and evidence $z$-score over $R$ (Eq.~\ref{eq:dz}) \\
$S_r$, $n_{\rm draw}$ & number of re-infill steps and of independent re-infill draws \\
$S$, $\rho$, $\mu,\sigma^2$, $\mu_c,\sigma^2_c$ & additive statistic $\sum_jI_j\Delta_j$; corrupted fraction; mean and variance of $\Delta_j$ on clean and on corrupted tokens (Eqs.~\ref{eq:mean}--\ref{eq:snr}) \\
$\phi_0$, $\phi_1$ & probability that a clean, respectively a corrupted, token is flagged in the signal-to-noise analysis; $\hat\phi_0,\hat\phi_1$ their observed values (called $\alpha$, $\beta$, $\alpha_{\rm obs}$, $\beta_{\rm obs}$ in the preregistration and the code) \\
$\gamma_0$, $\gamma_1$, $\gamma_{01}$ & correlation of $\Delta_j$ and $\Delta_k$, $j\neq k$, for two clean tokens, two corrupted tokens, and one of each (Proposition~\ref{prop:snrcorr}; theory only) \\
\bottomrule
\end{tabular}
\end{table}

\subsection{Masked scores}\label{ssec:scores}
For a token $j$ of modality $a$ with observed value $x_j$, four context levels are distinguished.
Let $\Gamma_{k(j)}$ be the group of $j$ in a partition $\Pi_K$, and let $x^a_{-\Gamma}$ denote the tokens of modality $a$ outside $\Gamma$.
\begin{align}
A_j &= \log p_\theta\big(x_j\mid x^a_{-\Gamma_{k(j)}},\,x^b\big), &
B_j &= \log p_\theta\big(x_j\mid x^a_{-\Gamma_{k(j)}},\,\varnothing^b\big),\notag\\
C_j &= \log p_\theta\big(x_j\mid \varnothing^a,\,x^b\big), &
D_j &= \log p_\theta\big(x_j\mid \varnothing^a,\,\varnothing^b\big).
\label{eq:levels}
\end{align}
$A_j$ uses the rest of the own modality and the other modality; $B_j$ uses the own modality only; $C_j$ uses the other
modality only; $D_j$ uses nothing and is the position marginal.
By subtraction, the total pointwise mutual information splits in two ways,
\begin{align}
A_j-D_j&=\underbrace{(B_j-D_j)}_{q_j:\ \text{intra-modal}}+\underbrace{(A_j-B_j)}_{\text{cross-modal given own context}}\notag\\
&=\underbrace{(A_j-C_j)}_{\text{intra-modal given }b}+\underbrace{(C_j-D_j)}_{\text{cross-modal}}.
\label{eq:pmi}
\end{align}
The quantity $q_j=B_j-D_j$ asks whether token $j$ is consistent with the rest of its own modality, relative to how likely
the value is at that position anyway.
It never sees the other modality, and therefore never sees the hypothesis that will be tested.
The cross-modal term $C_j-D_j$ is the decision evidence.
All four levels come from one primitive, Algorithm~\ref{alg:maskedscore}.

\begin{algorithm}[t]
\caption{\textsc{MaskedScore}: log-likelihood of target tokens under a mask pattern}\label{alg:maskedscore}
\Input{sequence $x$; target positions $P$; extra masked positions $X$; vocabulary slice $[lo,hi)$ of the targets}
\Output{$\lambda_j=\log p_\theta(x_j\mid\cdot)$ and $\pi^{\max}_j=\max_y p_\theta(y\mid\cdot)$ for $j\in P$}
\Fn{\textsc{MaskedScore}$(x,P,X,[lo,hi))$}{
  $\tilde x\leftarrow x$; set $\tilde x_P\leftarrow\mask$ and $\tilde x_X\leftarrow\mask$\;
  $\ell\leftarrow$ logits of $\tilde x$ \Comment*[r]{one network pass; batched with others}
  \For{$j\in P$}{
    $\eta_j\leftarrow\mathrm{logsoftmax}_{\fp}\big(\ell_{j,[lo,hi)}\big)$\;
    $\lambda_j\leftarrow\eta_j[x_j-lo]$ and $\pi^{\max}_j\leftarrow\exp\max\eta_j$\;
  }
  \Return $(\lambda,\pi^{\max})$\;
}
\end{algorithm}

\subsection{The gate statistic and the evidence-separation rule}\label{ssec:separation}
Two regimes are fixed per task before any run.
In the regime with \emph{no trusted reference} (detection, verification, matching), neither modality may be used to judge
the other, because the cross-modal relation is what is being tested.
In the regime with a \emph{trusted reference} (repair with a known clean caption), the trusted modality may be used in the gate,
as a predicted state is used in an innovation gate of tracking.
The paper studies the first regime.

\emph{Evidence-separation rule:} in the first regime, the gate may use only statistics computed with the other modality
masked, that is $B_j$ and $D_j$.
A gate that used $A_j$ with a candidate caption visible would give each candidate a different evidence set, and a gate that removed
low cross-modal-evidence caption words would remove exactly the words that make a wrong caption wrong.
The gate statistic is therefore
\begin{equation}
u_j=-q_j=-(B_j-D_j),
\label{eq:suspicion}
\end{equation}
where a larger $u_j$ means a more suspicious token.
The division by the position marginal $D_j$ is the optimal normalisation in one idealised case only: if a corrupted value is
drawn from the model's own position marginal independently of context, $q_j$ is the Neyman--Pearson statistic, whereas for
a uniform alternative, that is for corruptions that replace tokens by uniform random codes, the plain likelihood $B_j$ is.
Neither case is assumed to hold for JPEG, noise or splice, whose changed codes depend on the content, and no proposition about
detection power is claimed for them.
The null normalisation is therefore not assumed to improve detection, and the measurement does not show that it does.
% run: P2-GATE-M1
On the two uniform-code corruptions plain-$B$ has a slightly higher mean per-item AUROC than $q$ (patch 0.964 against 0.961,
swap 0.960 against 0.955), and on JPEG and noise the two differ by less than 0.01 (0.560 against 0.561, and 0.511 against 0.503;
no paired test was run on these two corruptions); Section~\ref{ssec:res_det} reports the full table.
Detection of content-dependent damage is an empirical question (Section~\ref{sec:protocol}).
The score $A_j-B_j$ without the null term cancels the own-context surprise of a corrupted token, because the term appears in
both $A_j$ and $B_j$; it is kept only as a control.

On the text side the gating unit is a word, that is a word-initial piece and its continuation pieces, and the suspicion of a
word is the maximum of the suspicions of its pieces (equivalently, the minimum of their $q$), following the idea that
multi-piece words need a word-level score~\cite{kauf2023plwl2r}.
Gating a word masks all its pieces.

\begin{algorithm}[t]
\caption{\textsc{ReliabilityScores}: suspicion of every token of one modality}\label{alg:reliability}
\Input{observed pair $\tilde x$; modality $a$; number of groups $K$; seed}
\Output{suspicion vector $u$ over the positions of $a$}
\Fn{\textsc{ReliabilityScores}$(\tilde x,a,K,\mathrm{seed})$}{
  draw a partition $\Pi_K$ of the positions of $a$ from the seed\;
  \For{$k=1,\dots,K$}{
    $B_{\Gamma_k}\leftarrow$ \textsc{MaskedScore}$(\tilde x,\ \Gamma_k,\ \text{positions of }b,\ \text{slice of }a)$ \Comment*[r]{$b$ withheld}
  }
  $D\leftarrow$ \textsc{MaskedScore}$(\tilde x,\ \text{all of }a,\ \text{all of }b,\ \text{slice of }a)$ \Comment*[r]{layout only; cached}
  $u_j\leftarrow-(B_j-D_j)$ for every position $j$ of $a$\;
  \lIf{$a=T$}{aggregate $u$ to words by the maximum over the pieces of a word}
  \Return $u$\;
}
\end{algorithm}

\subsection{Calibration by conformal risk control}\label{ssec:calibration}
The threshold is set on a calibration set $\mathcal D_{\rm cal}$ of $n_{\rm cal}$ clean pairs, disjoint by image from every other fold.
``Clean'' means images as distributed in the data set and passed through the tokenizer; corruption labels are relative to
that tokenisation.
For item $i$ with $m_i$ gating units and suspicions $u_{ij}$, the false-alarm proportion at threshold $\tau$ is
\begin{equation}
L_i(\tau)=\frac1{m_i}\sum_{j=1}^{m_i}\ind\big[u_{ij}>\tau\big]\in[0,1],
\label{eq:loss}
\end{equation}
which is non-increasing in $\tau$.
Define the empirical calibration risk $\widehat F_{\rm cal}(\tau)=n_{\rm cal}^{-1}\sum_i L_i(\tau)$ and the threshold
\begin{equation}
\hat\tau=\inf\Big\{\tau:\ \frac{n_{\rm cal}}{n_{\rm cal}+1}\,\widehat F_{\rm cal}(\tau)+\frac1{n_{\rm cal}+1}\le\alpha\Big\},
\label{eq:crc}
\end{equation}
with $\hat\tau=+\infty$ if the set is empty, so that nothing is flagged.
A token is flagged when $u_{ij}>\hat\tau$ with a strict inequality.

\begin{proposition}[calibration guarantee; conformal risk control applied]\label{prop:crc}
Let the calibration items and a new clean test item be exchangeable as items, and let $L_i$ be defined by
Eq.~\eqref{eq:loss}, with scores that have no ties.
Then $\E\big[L_{\rm test}(\hat\tau)\big]\le\alpha$.
\end{proposition}
\noindent\emph{Proof.} The loss is bounded by one and non-increasing in the threshold, so the statement is Theorem~1 of
conformal risk control~\cite{angelopoulos2024crc} with bound $B=1$. \qed

\paragraph{Wording and scope of the guarantee.}
The statement made in this paper, and no stronger one, is: the threshold is calibrated so that, for a clean test pair
exchangeable with the calibration pairs, the expected proportion of its tokens that are flagged is at most $\alpha$.
Three details keep the statement true in practice.
The bound is over items, so tokens inside an item need not be exchangeable or independent and positions may have different
score distributions.
Ties, which arise from \texttt{bf16} logits, are broken by an independent uniform jitter $\epsilon\sim U(0,\epsilon_0)$ added to
calibration and test scores, with $\epsilon_0$ below the smallest non-zero gap between distinct calibration scores, computed
and not chosen.
The statistic, $K$, $\alpha$ and the fold sizes are fixed before any test item is scored, because choosing them from calibration
or test items would break exchangeability.
The guarantee is marginal over the draw of the calibration set, so one realised calibration set may lie above $\alpha$ on a
finite test set.

\paragraph{What is not guaranteed.}
There is no bound on detection power: calibration controls false alarms on clean items only.
There is no bound on false alarms among the clean tokens of a \emph{corrupted} item, whose neighbours see contaminated context
and can be flagged more often; that rate is measured and reported (Section~\ref{sec:protocol}).
There is no guarantee under shift between calibration and test, such as another data set, caption style, resolution or
tokenizer.
Iterating the gate by re-scoring after masking breaks monotonicity in $\tau$, so the bound would not apply to it.
Words such as ``certified'' or ``provably reliable'' are not used for this gate.

\begin{algorithm}[t]
\caption{\textsc{CalibrateThreshold}}\label{alg:calibrate}
\Input{calibration set $\mathcal D_{\rm cal}$; modality $a$; level $\alpha$; $K$; seed rule}
\Data{clean pairs only, disjoint from all other folds}
\Output{threshold $\hat\tau$ and jitter scale $\epsilon_0$, stored with the hash of the item list, $K$, seed and statistic name}
\Fn{\textsc{CalibrateThreshold}$(\mathcal D_{\rm cal},a,\alpha,K)$}{
  \For{each item $i$ in $\mathcal D_{\rm cal}$}{
    $u_i\leftarrow$ \textsc{ReliabilityScores}$(x_i,a,K,\mathrm{seed}_i)$\;
  }
  $\epsilon_0\leftarrow$ a value below the smallest non-zero gap between distinct pooled scores\;
  add independent jitter $\epsilon_{ij}\sim U(0,\epsilon_0)$ to every $u_{ij}$\;
  evaluate $\widehat F_{\rm cal}(\tau)$ at every distinct pooled value and return the smallest $\tau$ satisfying Eq.~\eqref{eq:crc}\;
}
\end{algorithm}

\subsection{Gating}\label{ssec:gate}
The gate applies the frozen threshold to the suspicions of an observed pair.
Algorithm~\ref{alg:gate_set} computes the gated set $G$ and the reliable set $R$ with the same $K$, seed rule and statistic as
in calibration.
No iteration is used.

\begin{algorithm}[t]
\caption{\textsc{Gate} (uses Algorithms~\ref{alg:reliability} and~\ref{alg:calibrate})}\label{alg:gate_set}
\Input{observed pair $\tilde x$; modality $a$; frozen $\hat\tau$ and $\epsilon_0$; $K$}
\Output{gated set $G$, reliable set $R$, suspicion $u$}
\Fn{\textsc{Gate}$(\tilde x,a,\hat\tau,\epsilon_0,K)$}{
  $u\leftarrow$ \textsc{ReliabilityScores}$(\tilde x,a,K,\mathrm{seed})$\;
  add independent jitter $\epsilon_j\sim U(0,\epsilon_0)$ to every $u_j$\;
  $G\leftarrow\{j:u_j>\hat\tau\}$; $R\leftarrow$ complement of $G$\;
  \Return $(G,R,u)$\;
}
\end{algorithm}

\subsection{Withholding flagged tokens: GATE-M (marginalise)}\label{ssec:marginalise}
Masking a flagged token tells the network that its value is unknown.
For a network that equals the Bayes-optimal denoiser of the absorbing-state objective, masking the set $G$ marginalises it:
for a masked target $j\notin G$,
\begin{equation}
p_\theta(x_j\mid x_R,\,x_G=\mask)=\sum_{x_G}p(x_j\mid x_R,x_G)\,p(x_G\mid x_R)=p(x_j\mid x_R).
\label{eq:marginalise}
\end{equation}
This is the standard property of such denoisers~\cite{ou2025radd}, not a contribution.
A real network approximates it.
Structured patterns, such as a contiguous block of masked image tokens together with a fully masked caption, are rare under
random per-token masking, so the amortisation error may be larger there; a coherence test on the real model is part of the
protocol (Section~\ref{sec:protocol}).

When the set $G$ is selected from the observed values, Eq.~\eqref{eq:marginalise} is not enough.
Let $Y$ be the quantity to infer, for example the caption hypothesis.
Using $p(Y\mid\tilde v_R)$ equals the Bayes answer under the corruption model when four conditions hold:
(i) $R$ contains no corrupted token;
(ii) corrupted values carry no information about the clean values, which holds for replacement by uniform random codes
and fails for JPEG and noise;
(iii) corruption locations are independent of the content, which fails where changed codes concentrate in regions in which the
tokenizer is unstable;
(iv) no clean token is in $G$, or the clean tokens in $G$ are uninformative about $Y$.
The gate selects on the observed values, so the missingness is informative~\cite{rubin1976inference}.
Each violation has a cost: residual corrupted evidence in (i), discarded signal in (ii) because a noisy copy would call for a
denoising posterior and not for marginalisation, a wrong distribution $p(v_G\mid v_R)$ for selected positions in (iii), and lost
clean evidence at rate $\alpha$ in (iv).
The conditions therefore favour withholding for corruptions that satisfy (ii) and (iii) and are least favourable, possibly
not at all, for JPEG and noise; this is stated in advance and is not an account of any result.

\subsection{Image-direction evidence}\label{ssec:evidence}
The decision studied is paired verification: for one image, does the true caption $c^+$ explain the image better than a wrong
caption $c^-$.
The image-direction evidence uses one pass per caption with all image tokens masked and the caption visible.
Let $\Lambda^c_j[y]$ be the log-probability of code $y$ at image position $j$ in that pass, so that
\begin{equation}
\Delta_j=\Lambda^{c^+}_j[\tilde v_j]-\Lambda^{c^-}_j[\tilde v_j],\qquad
\bar\Delta_R=\frac1{|R|}\sum_{j\in R}\Delta_j,\qquad
Y_R=\ind\big[\bar\Delta_R>0\big].
\label{eq:delta}
\end{equation}
The position marginal $D_j$ cancels in the difference for a fixed image, so $\Delta_j=(C_j(c^+)-D_j)-(C_j(c^-)-D_j)$.
For a single candidate the per-caption evidence is $E_{I\leftarrow T}(c)=|R|^{-1}\sum_{j\in R}\big(\Lambda^{c}_j[\tilde v_j]-D_j[\tilde v_j]\big)$, and
$\bar\Delta_R=E_{I\leftarrow T}(c^+)-E_{I\leftarrow T}(c^-)$.
The logits of this pass do not depend on the image, which enters only through the gathered entries.
Hence all arms that differ in the set $R$ (no gate, gate, random gates, a changed-code gate, a denoised copy) reuse the same two
passes per item.
Ties and an empty set count as wrong, $Y=0$.
Two continuous summaries are used in the secondary analyses,
\begin{equation}
d_R=\frac{\bar\Delta_R}{s_R},\qquad z_R=\sqrt{|R|}\,d_R,
\label{eq:dz}
\end{equation}
where $s_R$ is the sample standard deviation of $\Delta_j$ over $R$.
The ratio $d_R$ rises when tokens with extreme per-token evidence are removed for any reason, which is why a random gate and
an uncertainty-ranked gate of equal size are reported next to it.

\begin{algorithm}[t]
\caption{\textsc{ImageEvidence} (GATE-M; uses Algorithms~\ref{alg:maskedscore} and~\ref{alg:gate_set})}\label{alg:evidence}
\Input{observed image tokens $\tilde v$; candidate captions $c_1,\dots,c_N$; reliable set $R$ from \textsc{Gate} or another arm; cached position marginal $D$}
\Output{$E_{I\leftarrow T}(c_k)$ for every candidate and, for a pair $(c^+,c^-)$, the paired quantities $Y_R$, $\bar\Delta_R$, $d_R$}
\Fn{\textsc{ImageEvidence}$(\tilde v,\{c_k\},R)$}{
  \For{each candidate $c_k$}{
    $\Lambda^{c_k}\leftarrow$ \textsc{MaskedScore} with all image tokens masked and $c_k$ visible, keeping the \fp{} log-probabilities over all image codes at all $m$ positions\Comment*[r]{one pass; reusable by every arm}
    $E_{I\leftarrow T}(c_k)\leftarrow\frac1{|R|}\sum_{j\in R}\big(\Lambda^{c_k}_j[\tilde v_j]-D_j[\tilde v_j]\big)$\;
  }
  \lIf{$N=2$}{$\Delta_j\leftarrow\Lambda^{c^+}_j[\tilde v_j]-\Lambda^{c^-}_j[\tilde v_j]$ for $j\in R$; \Return $\big(\ind[\bar\Delta_R>0],\ \bar\Delta_R,\ \bar\Delta_R/s_R\big)$}
  \Return $\{E_{I\leftarrow T}(c_k)\}$\;
}
\end{algorithm}

\paragraph{Analysis: when does withholding help an additive statistic?}\label{par:snr}
Consider the statistic $S=\sum_j I_j\Delta_j$ over $m$ tokens, with $I_j\in\{0,1\}$ the indicator that token $j$ is kept.
Under an independent-token approximation, clean tokens have mean $\mu$ and variance $\sigma^2$ and are kept with probability
$1-\phi_0$, corrupted tokens have mean $\mu_c$ and variance $\sigma_c^2$ and are kept with probability $1-\phi_1$, and a fraction
$\rho$ of tokens is corrupted.
Then
\begin{align}
\E S&=m\big[(1-\rho)(1-\phi_0)\mu+\rho(1-\phi_1)\mu_c\big],\label{eq:mean}\\
\Var S&=m\Big[(1-\rho)\big\{(1-\phi_0)\sigma^2+\phi_0(1-\phi_0)\mu^2\big\}\notag\\
&\qquad+\rho\big\{(1-\phi_1)\sigma_c^2+\phi_1(1-\phi_1)\mu_c^2\big\}\Big].\label{eq:var}
\end{align}
The no-gate case is $\phi_0=\phi_1=0$, and the gate raises $\E S/\sqrt{\Var S}$ if the ratio after exceeds the ratio before.
For $\mu_c=0$ this reduces to
\begin{equation}
\rho\,\sigma_c^2\,(\phi_1-2\phi_0+\phi_0^2)\;>\;\phi_0(1-\phi_0)(1-\rho)\,(\sigma^2+\mu^2).
\label{eq:snr}
\end{equation}
Three consequences are stated in advance.
A necessary condition is $\phi_1>2\phi_0-\phi_0^2$.
A random gate has $\phi_1=\phi_0$, for which the left side is negative, so it lowers the ratio.
On a clean item ($\rho=0$) any gate with $\phi_0>0$ lowers the ratio slightly.
The model is an approximation, since neighbouring tokens are correlated, and $\mu_c$ need not be zero even for random codes: a
caption that makes the predicted image distribution flatter gives random codes a higher probability.
The protocol estimates $\rho,\mu,\sigma^2,\mu_c,\sigma_c^2$ and the observed flag rates $\hat\phi_0,\hat\phi_1$ per corruption from
the per-token terms, and reports in how many corruptions the sign of the plug-in change of the ratio agrees with the sign of the
observed net accuracy difference (descriptive family E-pred, Section~\ref{ssec:families}); it is not a test and enters no claim
about the primary result.
E-pred uses Eqs.~\eqref{eq:mean}--\eqref{eq:var} as committed, and Proposition~\ref{prop:snrcorr} below is neither used by it nor
used to explain or interpret any E-pred result, agreeing or not.

\paragraph{Correlated tokens.}\label{par:snrcorr}
Eqs.~\eqref{eq:mean}--\eqref{eq:snr} treat the per-token terms as independent, which is an approximation for image tokens.
The following statement replaces independence by an exchangeable correlation within and between the two classes of tokens.
It is theory under stated assumptions and is not evaluated on the model.
The new symbols are $\gamma_0$, $\gamma_1$ and $\gamma_{01}$, the correlations of two clean, two corrupted, and one clean and one
corrupted token.

\begin{proposition}[withholding with exchangeably correlated tokens]\label{prop:snrcorr}
Condition on the corruption labels, so that $(1-\rho)m$ tokens are clean and $\rho m$ are corrupted, with the moments
$\mu,\sigma^2,\mu_c,\sigma_c^2$ and the keep probabilities $1-\phi_0$, $1-\phi_1$ of Eqs.~\eqref{eq:mean}--\eqref{eq:var}.
Assume (a) that for $j\neq k$ the correlation of $\Delta_j$ and $\Delta_k$ is $\gamma_0$ if both tokens are clean, $\gamma_1$ if both
are corrupted and $\gamma_{01}$ otherwise, with values for which the covariance matrix is positive semi-definite; and (b) that the
keep indicators $I_j$ are independent of each other and of every $\Delta_k$.
Then $\E S$ is given by Eq.~\eqref{eq:mean}, and, under (a) and (b),
\begin{equation}
\begin{aligned}
\Var S={}&V+(1-\rho)m\big((1-\rho)m-1\big)(1-\phi_0)^2\gamma_0\sigma^2\\
&+\rho m(\rho m-1)(1-\phi_1)^2\gamma_1\sigma_c^2\\
&+2\rho(1-\rho)m^2(1-\phi_0)(1-\phi_1)\gamma_{01}\sigma\sigma_c,
\end{aligned}
\label{eq:varcorr}
\end{equation}
where $V$ is the right side of Eq.~\eqref{eq:var}.
If moreover $\mu_c=0$ and $\E S>0$ with and without the gate, the gate raises $\E S/\sqrt{\Var S}$ if and only if
\begin{equation}
\begin{aligned}
&\rho\,\sigma_c^2\Big[(\phi_1-2\phi_0+\phi_0^2)\\
&\qquad+(\rho m-1)\,\gamma_1\,(\phi_1-\phi_0)(2-\phi_0-\phi_1)\Big]\\
&\quad+2\rho(1-\rho)\,m\,\gamma_{01}\,\sigma\sigma_c\,(1-\phi_0)(\phi_1-\phi_0)\\
&\qquad>\;(1-\rho)\,\phi_0(1-\phi_0)\,(\sigma^2+\mu^2).
\end{aligned}
\label{eq:snrcorr}
\end{equation}
Consequently:
(i) $\gamma_0$ does not enter Eq.~\eqref{eq:snrcorr}, and at $\gamma_1=\gamma_{01}=0$ it is Eq.~\eqref{eq:snr};
(ii) a gate with $\phi_1=\phi_0$, such as a random gate, lowers the ratio for every admissible correlation, and if $\gamma_{01}\ge0$ every
gate with $\phi_1\le\phi_0$ lowers it;
(iii) where $\phi_1>\phi_0$, a larger $\gamma_1$ or $\gamma_{01}$ can only turn a decrease of the ratio into an increase, so positive
correlation among corrupted tokens, or between clean and corrupted tokens, enlarges the region of $(\phi_0,\phi_1)$ in which the gate
raises the ratio, and as $(\rho m-1)\gamma_1\to\infty$ that region tends to $\{\phi_1>\phi_0\}$;
(iv) if $\gamma_{01}<0$, a gate that flags only clean tokens can raise the ratio;
(v) unlike Eq.~\eqref{eq:snr}, the condition depends on $m$.
\end{proposition}

\noindent\emph{Proof sketch.} Each of the $m(m-1)$ ordered pairs of distinct tokens adds $\E I_j\,\E I_k\,\mathrm{Cov}(\Delta_j,\Delta_k)$
to $\Var S$, which gives the three correlation terms of Eq.~\eqref{eq:varcorr}.
With $\mu_c=0$ the gate multiplies $\E S$ by $1-\phi_0$, so it raises the ratio if and only if
$(1-\phi_0)^2\,\Var S|_{\phi_0=\phi_1=0}>\Var S$; the clean-pair terms scale by $(1-\phi_0)^2$ on both sides and cancel, and the
remaining terms, divided by $m$, give Eq.~\eqref{eq:snrcorr}.
Appendix~\ref{app:snrcorr} gives the full proof. \qed

At zero correlation the statement is the independent-token analysis above: Eq.~\eqref{eq:varcorr} reduces to Eq.~\eqref{eq:var} and
Eq.~\eqref{eq:snrcorr} to Eq.~\eqref{eq:snr}.
Correlation therefore moves the gating frontier in the way (ii)--(iv) state, and independence is an approximation whose effect on the
frontier is of that kind and not of a fixed sign.
The ratio in Proposition~\ref{prop:snrcorr} is the signal-to-noise ratio of the sum.
If the per-token terms are jointly Gaussian, the accuracy $P(S>0)$ given the flags is $\Phi$ of the conditional ratio, which the
unconditional ratio approximates; without joint Gaussianity no statement about accuracy follows, and a gate can raise the ratio while
lowering the accuracy.
Assumption (b) is an approximation for GATE, whose flags are computed from the same token values as the evidence.
With correlation the ratio no longer grows like $\sqrt m$: it tends to a limit set by the shared component of the per-token terms.

\subsection{Text-direction evidence and conditioning arms}\label{ssec:text}
The text direction scores the caption tokens given the image.
With caption $c$ fully masked and the image visible, the evidence is the sum over caption tokens of the log-probability minus the same pass with the image
masked,
\begin{equation}
E_{T\leftarrow I}(c)=\sum_{j}\big(C_j-D_j\big),
\label{eq:text}
\end{equation}
where the image tokens in the gated set $G$ are masked in the first pass.
This is a conditioning arm in which the network, and not the arithmetic of a sum, withholds the flagged tokens.
In the exploratory analyses the statistic is the mean per-token value.

\subsection{Re-infilling flagged tokens: GATE-I (impute)}\label{ssec:reinfill}
Instead of leaving flagged tokens masked, the model can fill them in.
Algorithm~\ref{alg:reinfill} samples $n_{\rm draw}$ completions in $S_r$ steps, unmasking at step $s$ the next $k_s$ positions
of a random order drawn up front, and keeps the draw with the highest chain log-likelihood under the model itself, with the
caption masked (best of $n_{\rm draw}$; GATE-I).
The sampler is unguided and samples each position from the softmax at temperature one; it is the sampler fixed in the analysis plan
and not the sampler of the backbone's authors.
A hybrid variant that keeps only filled values whose recorded probability reaches $\kappa$ is described in the specification
and is not part of any preregistered family.
If the filled tokens are generated with the candidate caption visible, the evidence for that caption is biased upward.
Writing $\hat v_G\sim p(v_G\mid v_R,c)$, Bayes' rule gives
\begin{align}
\E\big[\log p(c\mid v_R,\hat v_G)\big]&=\log p(c\mid v_R)+\mathrm{KL}\big(p(v_G\mid v_R,c)\,\|\,p(v_G\mid v_R)\big)\notag\\
&\ge\log p(c\mid v_R),
\label{eq:leak}
\end{align}
so every candidate receives a non-negative bonus equal to how strongly it constrains the filled region, which need not favour the
true caption.
The rule is therefore to re-infill with the caption masked and to run the caption-visible variant only as a leakage control.
With exact sampling from $p(v_G\mid v_R)$, the expectation of $p(c\mid v_R,\hat v_G)$ equals $p(c\mid v_R)$, so GATE-I
is unbiased for GATE-M in the probability domain; in the log domain an $n_{\rm draw}$-sample estimate is biased downward by
Jensen's inequality and adds sampling variance.
Marginalising and imputing are not compared as a test: the preregistered re-infill family S2 compares the repaired flagged set
with repairing an equal-size random set (Section~\ref{ssec:families}).

\begin{algorithm}[t]
\caption{\textsc{ReInfill} (GATE-I; uses Algorithms~\ref{alg:maskedscore} and~\ref{alg:gate_set})}\label{alg:reinfill}
\Input{observed image tokens $\tilde v$; gated set $G$; steps $S_r$; draws $n_{\rm draw}$; seed}
\Output{the kept completed image $\hat v$ and the $n_{\rm draw}$ draws}
\Fn{\textsc{ReInfill}$(\tilde v,G,S_r,n_{\rm draw},\mathrm{seed})$}{
  \For{$\ell=1,\dots,n_{\rm draw}$}{
    $y\leftarrow\tilde v$ with $y_G\leftarrow\mask$; caption masked \Comment*[r]{caption masked: no leakage}
    draw a random order of $G$ from the seed; $\mathrm{chain}_\ell\leftarrow0$\;
    \For{$s=1,\dots,S_r$}{
      $U\leftarrow$ positions of $G$ still masked in $y$; $k_s\leftarrow\lceil|U|/(S_r-s+1)\rceil$\;
      $\pi_j\leftarrow p_\theta(\cdot\mid y)$ for $j\in U$ \Comment*[r]{one pass, via \textsc{MaskedScore}}
      take the next $k_s$ positions of $U$ in the drawn order; sample $y_j\sim\pi_j$; $\mathrm{chain}_\ell\mathrel{+}=\log\pi_j(y_j)$\;
    }
    $\hat v^{(\ell)}\leftarrow y$\;
  }
  \Return the draw $\hat v^{(\ell)}$ with the largest $\mathrm{chain}_\ell$ and the set of all draws\;
}
\end{algorithm}

\subsection{Two-way evidence fusion}\label{ssec:fusion}
For a candidate $c$ there are two directional statistics, $E_{T\leftarrow I}(c)$ of Eq.~\eqref{eq:text} and the image-direction
evidence of Eq.~\eqref{eq:delta}.
Inverse-variance weighting is not used, because the two statistics estimate different quantities on different scales, share
one network and are deterministic for a given input under full masking.
The fusion rule is reported in three steps on the paired differences $\Delta^{\rm pair}=(\Delta_T,\Delta_I)^\top$ with
$\Delta_d=E_d(c^+)-E_d(c^-)$:
(1) each direction alone;
(2) an unfitted Stouffer sum, with each paired margin standardised by its standard deviation on calibration items and the sum
divided by $\sqrt{2+2\hat r}$, where $\hat r$ is the calibration correlation;
(3) Fisher's linear discriminant,
\begin{equation}
\hat\omega=\hat\Sigma_{\Delta}^{-1}\hat\mu_{\Delta},\qquad F(c)=\hat\omega_T\,E_{T\leftarrow I}(c)+\hat\omega_I\,E_{I\leftarrow T}(c),
\label{eq:lda}
\end{equation}
with the weights fitted on pairs disjoint from the test pairs and frozen before the test fold.
The fusion therefore has two fitted weights, which is why the paper says ``no network finetuning'' and not ``training-free''.
The sum of per-token log-likelihood ratios in Eq.~\eqref{eq:delta} is the fusion rule of distributed detection~\cite{chair1986optimal}, which
equals the global log-likelihood ratio only for sources that are independent given the hypothesis.
Image tokens are not independent, so independence is an approximation here; Proposition~\ref{prop:snrcorr} states how correlation
moves the gating frontier, and it does not make the sum optimal.
The sum is used as a standard statistic and not as an optimal one.
Whether a fusion gain can be tested depends on the clean accuracy of the best single direction (Section~\ref{sec:protocol}).

\begin{algorithm}[t]
\caption{\textsc{GateVerify}: gated cross-modal verification (composes Algorithms~\ref{alg:maskedscore}--\ref{alg:reinfill})}\label{alg:gate}
\Input{observed pair $\tilde v$ and candidates $c_1,\dots,c_N$; frozen $\hat\tau_I,\epsilon_0$; fusion weights $\hat\omega$; arm $\in\{$none, GATE-M, GATE-I$\}$}
\Output{a score $F(c_k)$ per candidate and a decision}
\Fn{\textsc{GateVerify}$(\tilde v,\{c_k\},\hat\tau_I,\hat\omega,\text{arm})$}{
  $(G,R,u)\leftarrow$ \textsc{Gate}$(\tilde v,I,\hat\tau_I,\epsilon_0,K)$ \Comment*[r]{once per image; caption masked, hypothesis-free}
  \lIf{arm $=$ none}{$G\leftarrow\varnothing$; $R\leftarrow\{1,\dots,m\}$}
  \If{arm $=$ GATE-I}{$\hat v\leftarrow$ \textsc{ReInfill}$(\tilde v,G,S_r,n_{\rm draw},\mathrm{seed})$; $\tilde v\leftarrow\hat v$; $R\leftarrow\{1,\dots,m\}$}
  $\{E_{I\leftarrow T}(c_k)\}\leftarrow$ \textsc{ImageEvidence}$(\tilde v,\{c_k\},R)$\;
  \For{each candidate $c_k$}{
    $E_{T\leftarrow I}(c_k)\leftarrow$ \textsc{MaskedScore} on the caption, caption fully masked, image visible with the tokens of $G$ masked (GATE-M)\;
    $F(c_k)\leftarrow\hat\omega_T E_{T\leftarrow I}(c_k)+\hat\omega_I E_{I\leftarrow T}(c_k)$\;
  }
  \Return $\{F(c_k)\}$ and the decision (argmax, or a threshold for unpaired verification)\;
}
\end{algorithm}
